\documentclass[11pt,oneside]{article}

\pdfminorversion=7
\usepackage{QuizStyle}

\hypersetup{
    pdftitle={2026Fall Quiz for MATH301 Modern Algebra I},
    pdfauthor={Jungwoo Kim},
    pdfsubject={Quiz for MATH301 Modern Algebra I},
    pdfcreator={LaTeX}
}

\begin{document}

\quiztitle{2026Fall Quiz}

\begin{problem}{1.1}
Let $G$ be a finite group. Suppose that $G$ has no element of order $2$. Prove that $|G|$ is odd.
\end{problem}

\begin{solution}
For every $g\in G\setminus\{1_G\}$, we have $g\neq g^{-1}$; otherwise $g^2=1_G$, so $g$ would have order $2$. Hence the elements of $G\setminus\{1_G\}$ occur in disjoint pairs
\[
    \{g,g^{-1}\}.
\]
Thus $|G|-1$ is even, and therefore $|G|$ is odd.
\end{solution}

\begin{problem}{1.2}
Let $\sigma\in S_n$ have order $6$. Prove that if the cycle decomposition of $\sigma$ contains no $6$-cycle, then it must contain both a $2$-cycle and a $3$-cycle.
\end{problem}

\begin{solution}
Write the disjoint cycle decomposition of $\sigma$ as
\[
    \sigma=c_1c_2\cdots c_r,
\]
where the nontrivial cycles $c_i$ have lengths $m_i$. Then
\[
    |\sigma|=\lcmop(m_1,\ldots,m_r)=6.
\]
Hence each $m_i$ divides $6$. Since there is no $6$-cycle, every $m_i$ is $2$ or $3$. If there were no $2$-cycle, then $|\sigma|$ would divide $3$; if there were no $3$-cycle, then $|\sigma|$ would divide $2$. Both contradict $|\sigma|=6$. Therefore the cycle decomposition contains both a $2$-cycle and a $3$-cycle.
\end{solution}

\Needspace{24\baselineskip}
\begin{problem}{2.1}
Let $G$ be a group and let $\mathcal S$ be the set of all subgroups of $G$. Show that
\[
    g\cdot H=gHg^{-1},
\]
for $g\in G$ and $H\in\mathcal S$, defines a left group action of $G$ on $\mathcal S$.
\end{problem}

\begin{solution}
First, the proposed map is well defined. For $g\in G$ and $H\le G$, we have
\[
    1_G=g1_Gg^{-1}\in gHg^{-1}.
\]
If $gh_1g^{-1},gh_2g^{-1}\in gHg^{-1}$, then
\[
    (gh_1g^{-1})(gh_2g^{-1})^{-1}
    =g(h_1h_2^{-1})g^{-1}\in gHg^{-1},
\]
since $h_1h_2^{-1}\in H$. By the Subgroup Criterion, $gHg^{-1}\le G$.

Now let $g_1,g_2\in G$ and $H\in\mathcal S$. Then
\[
\begin{aligned}
    g_1\cdot(g_2\cdot H)
    &=g_1(g_2Hg_2^{-1})g_1^{-1}\\
    &=(g_1g_2)H(g_1g_2)^{-1}\\
    &=(g_1g_2)\cdot H,
\end{aligned}
\]
and
\[
    1_G\cdot H=1_GH1_G^{-1}=H.
\]
Thus $g\cdot H=gHg^{-1}$ defines a left group action of $G$ on $\mathcal S$.
\end{solution}

\begin{problem}{2.2}
Let $C_{12}=\langle x\rangle$ and $C_{18}=\langle y\rangle$ be cyclic groups of orders $12$ and $18$, respectively. Find, with proof, the number of homomorphisms
\[
    \varphi\colon C_{12}\to C_{18}.
\]
\end{problem}

\begin{notation}
Following MATH000, $C_n$ denotes an abstract cyclic group of order $n$. The original quiz uses $Z_n$ for the same object.
\end{notation}

\begin{solution}
A homomorphism $\varphi\colon C_{12}\to C_{18}$ is determined by the image of the generator $x$. Write
\[
    \varphi(x)=y^a,
    \qquad 0\le a<18.
\]
Since $x^{12}=1$, we must have
\[
    1=\varphi(x^{12})=y^{12a},
\]
so
\[
    18\mid12a
    \quad\Longleftrightarrow\quad
    3\mid a.
\]
Thus the only possible values are
\[
    a=0,3,6,9,12,15.
\]

Conversely, for each such $a$, define
\[
    \varphi_a(x^k)=y^{ak}.
\]
If $x^k=x^\ell$, then $12\mid(k-\ell)$; since $18\mid12a$, it follows that $18\mid a(k-\ell)$, so $y^{ak}=y^{a\ell}$. Hence $\varphi_a$ is well defined, and
\[
    \varphi_a(x^{k+\ell})
    =y^{a(k+\ell)}
    =y^{ak}y^{a\ell}
    =\varphi_a(x^k)\varphi_a(x^\ell),
\]
so $\varphi_a$ is a homomorphism. These six homomorphisms are distinct because they send $x$ to distinct elements of $C_{18}$. Therefore the number of homomorphisms is
\[
    \boxed{6}.
\]
\end{solution}


\Needspace{18\baselineskip}
\begin{problem}{3.1}
Let $G$ act on a nonempty set $A$. For $a\in A$, define the orbit of $a$ by
\[
    G\cdot a=\{g\cdot a\mid g\in G\}.
\]
Prove that if $b\in G\cdot a$, then $G\cdot b=G\cdot a$. Deduce that the distinct orbits form a partition of $A$.
\end{problem}

\begin{solution}
Suppose that $b\in G\cdot a$, so $b=g\cdot a$ for some $g\in G$. Then
\[
\begin{aligned}
    G\cdot b
    &=\{h\cdot(g\cdot a)\mid h\in G\}\\
    &=\{(hg)\cdot a\mid h\in G\}\\
    &=G\cdot a,
\end{aligned}
\]
since $\{hg\mid h\in G\}=G$.

Every $a\in A$ lies in $G\cdot a$, so the orbits cover $A$. If two orbits $G\cdot a$ and $G\cdot b$ intersect, choose $c\in(G\cdot a)\cap(G\cdot b)$. By the first part,
\[
    G\cdot a=G\cdot c=G\cdot b.
\]
Hence any two distinct orbits are disjoint, and the distinct orbits form a partition of $A$.
\end{solution}

\begin{problem}{3.2}
Let $G$ be a group and let $\varphi\in\Aut(G)$. Define
\[
    \operatorname{Fix}(\varphi)=\{g\in G\mid\varphi(g)=g\}.
\]
\begin{enumerate}[label=(\alph*),leftmargin=2.5em,itemsep=0.25em,topsep=0.4em]
    \item Prove that $\operatorname{Fix}(\varphi)$ is a subgroup of $G$.
    \item Let $G=\langle x\rangle$ be cyclic of order $12$, and let $\varphi$ be the automorphism determined by $\varphi(x)=x^5$. Determine $\operatorname{Fix}(\varphi)$.
\end{enumerate}
\end{problem}

\begin{solution}
\begin{enumerate}[label=(\alph*),leftmargin=2.5em,itemsep=0.45em,topsep=0.2em]
    \item Since $\varphi(1_G)=1_G$, the set $\operatorname{Fix}(\varphi)$ is nonempty. If $a,b\in\operatorname{Fix}(\varphi)$, then
    \[
        \varphi(ab^{-1})=\varphi(a)\varphi(b)^{-1}=ab^{-1}.
    \]
    Thus $ab^{-1}\in\operatorname{Fix}(\varphi)$, so $\operatorname{Fix}(\varphi)\le G$ by the Subgroup Criterion.

    \item Every element of $G$ has the form $x^k$, and
    \[
    \begin{aligned}
        x^k\in\operatorname{Fix}(\varphi)
        &\iff x^{5k}=x^k\\
        &\iff x^{4k}=1_G\\
        &\iff 12\mid4k\\
        &\iff 3\mid k.
    \end{aligned}
    \]
    Therefore
    \[
        \operatorname{Fix}(\varphi)
        =\langle x^3\rangle
        =\{1_G,x^3,x^6,x^9\}.
    \]
\end{enumerate}
\end{solution}

\end{document}
